% ============================================================
%  IEEE S&P 2026 Conference Poster
%  Title: Privacy-Preserving Malware Analysis with On-Prem LLMs
%         via Reinforcement Learning Alignment
%  Author: Gaspard Baye, Valix AI
%
%  Compile:  pdflatex poster_sp2026.tex   (twice for references)
%  Requires: beamer, beamerposter, tikz, booktabs, amsmath
%
%  Print size: 24 x 36 inches (portrait).
%  Scale to US letter PDF before uploading to HotCRP.
% ============================================================

\documentclass[final]{beamer}

\usepackage[
  orientation=portrait,
  size=custom,
  width=60.96,    % 24 inches in cm
  height=91.44,   % 36 inches in cm
  scale=1.15
]{beamerposter}

\usepackage{amsmath,amssymb}
\usepackage{booktabs}
\usepackage{graphicx}
\usepackage{tikz}
\usepackage{url}
\usepackage{xcolor}
\usepackage[nocompress]{cite}

\usetikzlibrary{arrows.meta, shapes.geometric, positioning}

% ============================================================
%  COLOUR PALETTE
% ============================================================
\definecolor{spnavy}{HTML}{0D1B2A}
\definecolor{spaccent}{HTML}{F4A024}
\definecolor{spblue}{HTML}{1B6CA8}
\definecolor{splightblue}{HTML}{D6E8F7}
\definecolor{spgray}{HTML}{F5F5F5}
\definecolor{spred}{HTML}{C0392B}
\definecolor{spgreen}{HTML}{27AE60}

% ============================================================
%  BEAMER THEME SETUP
% ============================================================
\usetheme{default}
\setbeamercolor{background canvas}{bg=white}
\setbeamercolor{block title}{fg=white, bg=spblue}
\setbeamercolor{block body}{fg=black, bg=splightblue}
\setbeamercolor{block title alerted}{fg=white, bg=spred}
\setbeamercolor{block body alerted}{fg=black, bg=spgray}
\setbeamercolor{block title example}{fg=white, bg=spgreen}
\setbeamercolor{block body example}{fg=black, bg=spgray}
\setbeamertemplate{navigation symbols}{}
\setbeamerfont{block title}{size=\large, series=\bfseries}
\setbeamerfont{block body}{size=\normalsize}

\newcommand{\vsep}{\vspace{1.2em}}

% ============================================================
\begin{document}
\begin{frame}[t]{}

% ============================================================
%  HEADER — unified single box: logo | title | badge
% ============================================================
\begin{beamercolorbox}[wd=\paperwidth, ht=14cm, dp=1cm,
    leftskip=1cm, rightskip=1cm, sep=0pt]{block title}

  \vspace{0.6cm}

  % --- Three-column layout: logo | title | conference badge ---
  \begin{minipage}[c]{0.16\paperwidth}
    % Replace with \includegraphics[height=3cm]{valix_logo} when available
    {\fontsize{26}{30}\selectfont\bfseries\color{spaccent} Valix\,AI}\\[0.3em]
    {\small\color{white} Chief Technology Officer}
  \end{minipage}
  \hfill
  \begin{minipage}[c]{0.60\paperwidth}
    \centering
    {\fontsize{30}{38}\selectfont\bfseries\color{white}
      Privacy-Preserving Malware Analysis\\[0.2em]
      with On-Prem LLMs via\\[0.2em]
      Reinforcement Learning Alignment}
  \end{minipage}
  \hfill
  \begin{minipage}[c]{0.16\paperwidth}
    \raggedleft
    {\large\bfseries\color{spaccent} IEEE S\&P 2026}\\[0.3em]
    {\normalsize\color{white} Poster Session}\\[0.3em]
    {\small\color{white} San Francisco, CA}
  \end{minipage}

  \vspace{0.6cm}
  {\color{spaccent}\rule{\linewidth}{1.5pt}}
  \vspace{0.3cm}

  % --- Author line centred below ---
  \centering
  {\Large\bfseries\color{white} Gaspard Baye}\quad
  {\normalsize\color{white}
    Chief Technology Officer, Valix AI\quad
    \texttt{baye@valix.ai}}

  \vspace{0.4cm}
\end{beamercolorbox}

% ============================================================
%  TWO-COLUMN BODY
% ============================================================
\begin{columns}[T]

% ============================================================
%  LEFT COLUMN
% ============================================================
\begin{column}{0.48\linewidth}

% ------ MOTIVATION ------------------------------------------
\begin{block}{Motivation}
SOC analysts must classify malware behaviors from sandbox reports,
map them to MITRE ATT\&CK tactics, and prioritize responses under high
alert volume.
Frontier LLMs (GPT-4, Gemini) automate this triage well, but
\textbf{require sensitive endpoint telemetry to leave the organization},
creating conflicts with SOC~2, HIPAA, and NIST~800-53.

\vsep

% GRAPHIC 1 — Cloud vs On-Prem diagram
% Two side-by-side flow diagrams separated by a firewall line.
% Left (red): SOC Analyst -> Cloud API, labelled "Telemetry leaves perimeter"
% Right (green): SOC Analyst -> On-Prem GPU, labelled "Data stays inside firewall"
\begin{center}
\begin{tikzpicture}[
  flowbox/.style={draw, rounded corners=4pt,
                  minimum width=2.8cm, minimum height=1cm,
                  align=center, font=\small},
  flowarr/.style={-{Stealth[length=5pt]}, thick}
]
  % Left panel: cloud risk
  \node[flowbox, fill=red!10, draw=spred]   (soc1)  at (0,0)    {SOC\\Analyst};
  \node[flowbox, fill=red!10, draw=spred]   (cloud) at (5,0)    {Cloud\\API};
  \draw[flowarr, color=spred, dashed]
    (soc1) -- (cloud)
    node[midway, above, font=\tiny, color=spred]{sensitive telemetry};
  \node[font=\small\bfseries, color=spred] at (2.5,-1.4) {Privacy Risk};

  % Firewall divider
  \draw[dashed, thick, color=spblue] (7.2,-1.8) -- (7.2,1.2)
    node[above, font=\footnotesize\color{spblue}]{Firewall};

  % Right panel: on-prem safe
  \node[flowbox, fill=green!10, draw=spgreen] (soc2)   at (9.5,0)  {SOC\\Analyst};
  \node[flowbox, fill=green!10, draw=spgreen] (onprem) at (14.5,0) {On-Prem\\LLM};
  \draw[flowarr, color=spgreen]
    (soc2) -- (onprem)
    node[midway, above, font=\tiny, color=spgreen]{local inference};
  \node[font=\small\bfseries, color=spgreen] at (12,-1.4) {Privacy Safe};
\end{tikzpicture}
\end{center}

\vsep

\begin{alertblock}{Core Research Question}
  \centering\large
  Can a \textbf{compact (3B) on-prem LLM}, aligned via RL, approach
  frontier accuracy for malware triage -- with \textbf{zero data exfiltration}?
\end{alertblock}

\end{block}

\vsep

% ------ THREAT MODEL ----------------------------------------
\begin{block}{Threat Model}
\begin{itemize}
  \item \textbf{Adversary:} Network-level attacker or cloud-provider
    insider able to intercept or correlate queries to a cloud API.
  \item \textbf{Protected asset:} Behavioral telemetry of the
    organization's endpoints, network flows, and malware sandbox data.
  \item \textbf{Security property:} On-prem inference confines all
    processing inside the security perimeter; no query crosses a
    boundary, eliminating the data-exfiltration attack surface.
\end{itemize}
\end{block}

\vsep

% ------ DISTILLATION ----------------------------------------
\begin{block}{Silver-Standard Distillation (Step 0)}

\begin{columns}[T]
  \begin{column}{0.54\linewidth}
    We curate \textbf{10,000 raw malware sandbox reports} from
    Hybrid-Analysis and VirusTotal.
    Gemini~2.0-Flash generates reasoning chains, which are then
    semantically augmented to form a \textbf{1M-sample HuggingFace dataset}:
    \[
      \mathcal{D} = \{(x_i,\, c_i,\, y_i)\}_{i=1}^{1,000,000}
    \]
    where $x_i$ is the raw report, $c_i$ the reasoning chain,
    $y_i$ the threat label.

    \smallskip
    \textit{Limitation:} Gemini errors can propagate into $\mathcal{D}$.
    Rule-based RL rewards in Stages~2--3 act as a correction filter.
  \end{column}
  \begin{column}{0.40\linewidth}
    % GRAPHIC 2 — Vertical distillation flow
    % Raw Report -> Gemini 2.0-Flash -> <thought> chain -> Corpus D
    \begin{center}
    \begin{tikzpicture}[
      fbox/.style={draw, rounded corners=3pt, fill=spgray,
                   minimum width=3.2cm, minimum height=0.85cm,
                   align=center, font=\small},
      farr/.style={-{Stealth}, thick, color=spblue}
    ]
      \node[fbox] (rpt) at (0,0)    {Raw Reports (10k)};
      \node[fbox] (gem) at (0,-2.2) {Gemini 2.0-Flash};
      \node[fbox] (chn) at (0,-4.4) {Semantic Augmentation};
      \node[fbox] (cor) at (0,-6.6) {1M-Sample HF Dataset};
      \draw[farr] (rpt) -- (gem);
      \draw[farr] (gem) -- (chn);
      \draw[farr] (chn) -- (cor);
    \end{tikzpicture}
    \end{center}
  \end{column}
\end{columns}

\end{block}

\vsep

% ------ CONTRIBUTIONS ----------------------------------------
\begin{block}{Key Contributions}
\begin{enumerate}
  \item \textbf{First} ATT\&CK-grounded multi-objective reward design
    applied to RL alignment for SOC reasoning.
  \item \textbf{Silver-Standard Distillation:} frontier-model reasoning
    bootstrapped into a local corpus without human annotation.
  \item \textbf{Privacy-accuracy tradeoff analysis} between cloud and
    on-premise LLMs in a malware triage setting.
\end{enumerate}
\end{block}

\end{column} % end left column

% ============================================================
%  RIGHT COLUMN
% ============================================================
\begin{column}{0.48\linewidth}

% ------ PIPELINE --------------------------------------------
\begin{block}{Three-Stage RL Alignment Pipeline}

% GRAPHIC 3 — Horizontal pipeline flow diagram
% [Corpus D] -> [Stage 1: SFT] -> [Stage 2: PPO] -> [Stage 3: GRPO] -> [On-Prem Expert]
% Each box colour-coded; annotation below each stage.
\begin{center}
\begin{tikzpicture}[
  pbox/.style={draw, rounded corners=5pt,
               minimum width=3cm, minimum height=2cm,
               align=center, font=\small\bfseries},
  parr/.style={-{Stealth[length=7pt]}, line width=1.5pt, color=spnavy},
  plbl/.style={font=\tiny, text width=3cm, align=center, color=gray}
]
  \node[pbox, fill=spblue!20,    draw=spblue]   (sft)  at (0,0)
    {Stage 1\\SFT\\{\footnotesize Learn format}};
  \node[pbox, fill=spaccent!25,  draw=spaccent] (ppo)  at (5.2,0)
    {Stage 2\\PPO\\{\footnotesize Reward quality}};
  \node[pbox, fill=spgreen!20,   draw=spgreen]  (grpo) at (10.4,0)
    {Stage 3\\GRPO\\{\footnotesize Rank \& refine}};
  \node[pbox, fill=spnavy!15,    draw=spnavy]   (out)  at (15.6,0)
    {On-Prem\\SOC Expert};

  \draw[parr] (sft)  -- (ppo);
  \draw[parr] (ppo)  -- (grpo);
  \draw[parr] (grpo) -- (out);

  \node[plbl] at (0,-1.8)   {Imitates Gemini reasoning};
  \node[plbl] at (5.2,-1.8) {ATT\&CK reward signal};
  \node[plbl] at (10.4,-1.8){Group-relative advantage};
  \node[plbl] at (15.6,-1.8){On-prem deployment};
\end{tikzpicture}
\end{center}

\vsep

\textbf{Stage~1 -- SFT.}
Model imitates distilled \texttt{<thought>}/\texttt{<answer>} structure:
$\;\mathcal{L}_\text{SFT}(\theta) = -\sum_i \log P_\theta(y_i, c_i \mid x_i)$.

\vsep
\textbf{Stage~2 -- PPO.}
A composite reward is maximised while a frozen reference model
$\pi_\text{ref}$ (KL anchor) prevents reward exploitation.

\vsep
\textbf{Stage~3 -- GRPO.}
$G{=}8$ candidate responses per report; within-group advantage
$A_k {=} (r_k {-} \mu(r))\,/\,\sigma(r)$ drives the policy toward
the best response without a separate critic network.

\end{block}

\vsep

% ------ REWARD ----------------------------------------------
\begin{exampleblock}{Multi-Objective Reward Architecture}

Each response scored on three analyst-relevant criteria:
\[
  R(O \mid Q)=
    \underbrace{\lambda_\text{fmt}\,R_\text{fmt}}_{\text{structured output}}
  +\underbrace{\lambda_\text{acc}\,R_\text{acc}}_{\text{correct answer}}
  +\underbrace{\lambda_\text{prec}\,R_\text{prec}}_{\text{ATT\&CK terms}}
\]

\begin{center}
\begin{tabular}{lll}
  \toprule
  \textbf{Component} & \textbf{Range} & \textbf{Rewards} \\
  \midrule
  $R_\text{fmt}$  & $[0,\,0.5]$ & \texttt{<thought>}/\texttt{<answer>} tags present \\
  $R_\text{acc}$  & $\{0,\,1\}$ & Exact ground-truth threat label match \\
  $R_\text{prec}$ & $[0,\,1]$   & Density of MITRE ATT\&CK terms in response \\
  \bottomrule
\end{tabular}
\end{center}

\smallskip
% GRAPHIC 4 — Reward component bar chart (generate externally)
% Grouped horizontal bars: R_fmt, R_acc, R_prec for each model.
% All baselines show 0 on R_fmt (no tags). Show expected post-GRPO uplift.
% Recommended: pgfplots or matplotlib PDF, then \includegraphics{reward_bars}
\textit{(Reward component breakdown --- full results at poster session.)}

\end{exampleblock}

\vsep

% ------ RESULTS ---------------------------------------------
\begin{block}{Preliminary Results \textnormal{(Pilot, $N\!=\!20$)}}

\small\textbf{Note:} Pilot benchmark --- interpret as feasibility signal.
Full statistical evaluation with confidence intervals to follow.

\vsep

\begin{center}
\begin{tabular}{lcccc}
  \toprule
  \textbf{Model} & \textbf{Acc.} & \textbf{Tech.\ Prec.}
                 & \textbf{Words} & \textbf{Reasoning} \\
  \midrule
  Gemini 2.0 (cloud ref.)   & 98\% & 0.108 & 505 & 0.00 \\
  LLaMA-3 8B (baseline)     & 85\% & 0.082 & 208 & 0.00 \\
  Mistral 7B (baseline)     & 75\% & 0.091 & 282 & 0.00 \\
  \midrule
  SFT Expert (LLaMA-3.2)   & \multicolumn{4}{c}{\emph{poster session}} \\
  PPO Expert (LLaMA-3.2)   & \multicolumn{4}{c}{\emph{poster session}} \\
  GRPO Expert (LLaMA-3.2)  & \multicolumn{4}{c}{\emph{poster session}} \\
  \bottomrule
\end{tabular}
\end{center}

\vsep

% GRAPHIC 5 — Accuracy vs. Word Count scatter
% X-axis: Avg words per response. Y-axis: Accuracy.
% Points: Gemini (star), LLaMA-3 (circle), Mistral (triangle).
% Shade a "target zone" (high accuracy, low word count).
% Recommended: matplotlib PDF -> \includegraphics{scatter_acc_words}
\textit{(Accuracy vs.\ verbosity scatter --- full results at poster session.)}

\vsep
\textbf{Key insight:} No baseline produces structured reasoning
(Reasoning Density\,$=\!0$).
Most common failure: \emph{partial label extraction} -- correct malware
family identified, secondary behavior missed (e.g., persistence mechanism).
This is exactly the gap the \texttt{<thought>} chain and GRPO reward target.

\end{block}

\vsep

% ------ REPRODUCIBILITY -------------------------------------
\begin{block}{Implementation Details \& Reproducibility}

\begin{tabular}{ll}
  \textbf{Base model:}   & LLaMA 3.2 (3B) via HuggingFace TRL \\
  \textbf{Hardware:}     & 1$\times$ NVIDIA H100 80\,GB \\
  \textbf{Dataset:}      & 1M-sample HF dataset (augmented from 10k reports) \\
  \textbf{Distillation:} & Gemini 2.0-Flash (structured TTP reasoning) \\
  \textbf{SFT:}          & 1 epoch, lr\,$=2{\times}10^{-4}$, cosine schedule \\
  \textbf{PPO/GRPO:}     & lr\,$=5{\times}10^{-6}$, group size $G{=}8$ \\
  \textbf{Evaluation:}   & 20 human-verified malware analysis questions \\
\end{tabular}

\vsep
\textbf{Code, prompts, and reward scripts will be released upon
acceptance} to support community reproducibility.

\end{block}

\vsep

% ------ REFERENCES ------------------------------------------
{\small
\begin{thebibliography}{6}
\bibitem{ppo}    Schulman et al., PPO, arXiv:1707.06347, 2017.
\bibitem{grpo}   Shao et al., DeepSeekMath, arXiv:2402.03300, 2024.
\bibitem{llama3} Dubey et al., Llama~3, arXiv:2407.21783, 2024.
\bibitem{mitre}  MITRE, ATT\&CK, \url{https://attack.mitre.org}, 2015.
\bibitem{trl}    von Werra et al., TRL, \url{https://github.com/huggingface/trl}.
\bibitem{ctibench} Ameri et al., CTIBench, arXiv:2406.07521, 2024.
\end{thebibliography}
}

\end{column} % end right column
\end{columns}

\vsep

% ============================================================
%  FOOTER — plain coloured box, no overlay
% ============================================================
\begin{beamercolorbox}[wd=\paperwidth, ht=2.5cm, dp=0pt,
    leftskip=1.5cm, rightskip=1.5cm]{block title}
  \vfill
  {\small\color{white}
    IEEE Symposium on Security and Privacy 2026 \quad$|$\quad
    Valix AI \quad$|$\quad \texttt{baye@valix.ai}}
  \hfill
  {\small\bfseries\color{spaccent} Poster Session}
  \vfill
\end{beamercolorbox}

\end{frame}
\end{document}
